\section{Conclusiones}

Se construyó en \matlab{} una aplicación con interfaz gráfica que integra los métodos numéricos del curso para analizar la vibración real de un automóvil medida con el acelerómetro de un teléfono. La aplicación calcula la duración de la señal, detecta e interpola datos faltantes, obtiene la derivada y la integral, resuelve la ecuación diferencial del movimiento y ajusta líneas de tendencia por mínimos cuadrados. Ofrece menús para elegir algoritmos y parámetros, gráficas de los datos originales y procesados, indicadores numéricos, reproducción en tiempo real con grabación controlada por el usuario y exportación con nombre de archivo definido por él. Con esto se cubren todos los requisitos del proyecto.

Los métodos se programaron sin usar las funciones equivalentes de \matlab{} y se validaron contra una solución analítica. Los órdenes de convergencia observados coinciden con los teóricos, y la diferencia frente a \texttt{interp1}, \texttt{gradient}, \texttt{trapz} y \texttt{polyfit} es del orden del redondeo de máquina cuando el método es equivalente.

Con las mediciones reales, la aplicación identificó la frecuencia de encendido del motor (25.26~Hz, unas 758~rpm con el motor frío) y registró su descenso a 22.86~Hz (686~rpm) al calentarse, en concordancia con la medición independiente de phyphox. También mostró que integrar la aceleración sin corregir el sesgo produce un desplazamiento falso de más de 100~m en un vehículo estacionado y que la tendencia por mínimos cuadrados lo reduce a centímetros.

La principal lección del análisis es que el método más efectivo depende de la señal. Con una señal suave y bien muestreada se cumple la jerarquía teórica: trazador antes que lineal, cinco puntos antes que centrada, Simpson antes que trapecio y RK4 antes que Euler. Con la vibración real, a unas cuatro muestras por ciclo, la interpolación lineal superó a las cúbicas, y el error de integración lo dominó el sesgo del sensor, no el orden del integrador. Elegir un algoritmo de forma justificada exige comparar los métodos con los datos del problema, y la aplicación incluye una pestaña para hacerlo.

Como trabajo futuro se proponen tres mejoras: medir simultáneamente en el cofre y en la cabina para cuantificar el aislamiento de los soportes del motor; usar un sensor externo con mayor frecuencia de muestreo para observar los armónicos superiores; y fusionar el acelerómetro con el giroscopio para estimar la velocidad del vehículo independientemente de la orientación del teléfono.
